% Expository adaptation, October 7, 2026. Original mathematics: OpenAI, September 23, 2026.
% Prose revised and explanations expanded; see README.md for the source and scope.
\section{Conditional averaging in the shared environment}
\label{sec:averaging}

The random tables have the right conditional means. The estimator cannot average over every possible environment at every recursive call, so it needs a shorter averaging operation. Moreover, ordinary averaging is not enough: the queried source and destination must stay in their prescribed retained sets.

Lemma~\ref{lem:env-gap} is the guarantee used by the estimator. A walk of length $d$ reduces the discrepancy from the appropriate conditional mean exponentially in $d$, with constants independent of the field size and retention rates. Lemma~\ref{lem:env-variants} gives a second, computational guarantee: a path through the averaging walk can be described without saving the vertex identifiers whose retention is being conditioned on. The tests involving an identifier can be postponed until the algorithm visits that vertex.

We first prove the mixing guarantee from a uniform spectral gap, and then explain the identifier-independent path description.

The identifiers are fixed when the conditional space is defined.  We
condition on one retained identifier per conditioned channel; additional
independent bit coordinates may be varied or frozen as specified below.

The \emph{full environment law} already includes the frame restriction
on each channel.  A conditional slice additionally imposes the indicated
retention events and fixes any designated frozen bits.  Its averaging
operator and $L_s$ norms use the normalized probability law on that slice.

The group-theoretic input ensures that a fixed list of transformations mixes uniformly over an entire family of finite spaces. In the representation-theoretic formulation, a vector that changes little under every generator must be close to a vector invariant under the whole group. On a finite transitive action the invariant functions are constants, so this gives a quantitative gap between constants and all mean-zero functions.

The only external group-theoretic inputs are property~(T) of
$\mathrm{SL}_H(\mathbb Z)$ for $H\ge3$ and relative property~(T) of
\[
(\mathrm{SL}_2(\mathbb Z)\ltimes\mathbb Z^2,\mathbb Z^2).
\]
We use Shalom's quantitative statements:
\cite[Theorem~2.6, p.~156]{Shalom1999} for the linear group and
\cite[Theorem~2.1, p.~152]{Shalom1999} for the relative property.
We derive the affine finite-action gap and the conditional sampler from
these inputs below.

\begin{lemma}[A uniform affine spectral gap]\label{lem:env-affine}
Let
$G_H=\mathrm{SL}_H(\mathbb Z)\ltimes\mathbb Z^H$, with fixed
$H\ge3$. There is a fixed finite symmetric generating multiset and a
lazy rational random walk on it whose action on every finite transitive
$G_H$-set has a spectral gap bounded below by a positive constant
depending only on $H$.
\end{lemma}

\begin{proof}
We first control approximate invariance under translations, then under the linear group, and finally convert this control into a mixing bound for finite actions. Let $\mathcal S$ consist of the elementary matrices with a single
$\pm1$ off the diagonal, acting linearly, and the translations by
$\pm e_i$, where $e_i$ are the standard basis vectors of $\mathbb Z^H$.
This is a fixed finite symmetric generating set of $G_H$.
For a unitary representation $\rho$ and a vector $\xi$, put
\[
\delta=\max_{g\in\mathcal S}\|\rho(g)\xi-\xi\|.
\]
We first bound the distance from $\xi$ to the invariant subspace.

For each coordinate pair $\{i,j\}$, embed the two-dimensional affine
group on those coordinates, and let $P_{ij}$ be the orthogonal projection
onto the subspace fixed by its translations. This subspace and its
orthogonal complement are invariant under that pair group, because its
translation subgroup is normal. The eight generators in
\cite[Theorem~2.1, p.~152]{Shalom1999} belong to $\mathcal S$. Applying their
relative Kazhdan constant $1/10$ on the orthogonal complement gives
\[
\|\xi-P_{ij}\xi\|\le10\delta.
\]
This is the projection argument in the proof of
\cite[Corollary~2.3, p.~154]{Shalom1999}.

To combine the two-coordinate estimates, project successively onto their translation-invariant subspaces. The projections $P_{ij}$ commute: they are limits of averages of
commuting coordinate translations. Their product $P_{\rm tr}$ projects onto
the subspace fixed by all translations. With $q=\binom H2$, telescoping
the product gives
\[
\|\xi-P_{\rm tr}\xi\|
\le\sum_{i<j}\|\xi-P_{ij}\xi\|
\le10q\delta.
\]
Since the full translation subgroup is normal, $P_{\rm tr}$ commutes
with every $\rho(g)$. Its range carries a representation of
$G_H/\mathbb Z^H=\mathrm{SL}_H(\mathbb Z)$. Let $P_{\rm inv}$ project onto the
full $G_H$-invariant subspace, and let $\kappa_H>0$ be the Kazhdan
constant for the elementary generators from
\cite[Theorem~2.6, p.~156]{Shalom1999}; adjoining inverse generators
preserves the bound. Applying that constant on the range of
$P_{\rm tr}$ gives
\[
\|P_{\rm tr}\xi-P_{\rm inv}\xi\|\le\kappa_H^{-1}\delta,
\qquad
\|\xi-P_{\rm inv}\xi\|\le C_H\delta,
\quad C_H=10q+\kappa_H^{-1}.
\]
Thus the same generating set controls distance to invariants in every
unitary representation, with a constant depending only on $H$.

We now apply the representation estimate to the finite space on which the averaging algorithm will run. Let $\rho$ be the permutation representation of a finite transitive
action, with the uniform probability measure, and put $m=|\mathcal S|$.
The rational Markov operator
\[
K_H=\frac12I+\frac1{2m}\sum_{g\in\mathcal S}\rho(g)
\]
is self-adjoint, and its spectrum lies in $[0,1]$. The invariant
functions are exactly the constants. For $f$ orthogonal to constants,
the preceding inequality and the Dirichlet identity give
\[
\langle f,(I-K_H)f\rangle
=\frac1{4m}\sum_{g\in\mathcal S}\|\rho(g)f-f\|_2^2
\ge\frac1{4mC_H^2}\|f\|_2^2.
\]
This proves the uniform spectral gap.
\end{proof}

The conditioning event has a particularly useful coordinate description. Once the nonconstant coefficients and the value at the conditioned identifier are specified, the constant coefficient is determined. Thus the conditional space separates into a full-rank coefficient array and a uniform point in a box. The affine action can mix these coordinates even though their moduli differ.

\begin{lemma}[Conditional averaging]\label{lem:env-gap}
Consider the fixed product of rank channels, conditioned on retention of
one specified identifier per conditioned channel. Include any designated
set of independent uniform bits, with length padded to a multiple of
$H$, and freeze any remaining independent bits.
There is a rational Markov averaging operator $K$ on this conditional
space such that its stationary distribution is the indicated uniform
conditional law and, for each fixed $2\le s<\infty$,
\[
\|K^d-\Pi\|_{L_s\to L_s}\le C_s2^{-c_sd}
\qquad(d\ge0).
\]
Here $\Pi$ is conditional expectation over the variables being averaged.
The constants are independent of the field size, retention side
lengths, bit lengths, indicated identifiers, and frozen bits.
\end{lemma}

\begin{proof}
On a channel conditioned at $y$ with side $L$, use coordinates
$(F,w)$, where $F$ is its field frame and
$w=Av(y)\in\{0,\ldots,L-1\}^H$. The constant column is uniquely
determined by $F,w,y$. As observed above, $F$ is uniform over frames
and $w$ is independent uniform in its box. Identify this box with
$(\mathbb Z/L\mathbb Z)^H$.
For $(M,c)\in G_H$ act by
\[
(F,w)\longmapsto(MF,Mw+c),
\]
using reduction modulo $P$ on the first coordinate and modulo $L$
on the second.

This action is transitive. Reduction of $\mathrm{SL}_H(\mathbb Z)$
contains the elementary generators of $\mathrm{SL}_H(\mathbb F_P)$.
That finite group acts transitively on ordered frames of length
$r<H$: extend either frame to a basis and adjust a complementary
column to make the change-of-basis determinant one. Once the frame
has been moved, a translation adjusts $w$ to any specified value.
The two moduli need not agree. On an unconditioned channel use a
public identifier and side $P$, giving the same description.

If the free bit length is $Hq$, identify its strings with
$(\mathbb Z/2^q\mathbb Z)^H$ and use the affine action on that
coordinate alone. It is transitive. All currently free bit blocks
can be concatenated into this one component; frozen bits are not
transformed.
There are a bounded number of components. Apply the lazy walk from
Lemma~\ref{lem:env-affine} to a uniformly selected component.
If there are $J_{\rm comp}$ components and the single-component gap is
at least $\gamma_H$, the product walk has gap at least
$\gamma_H/J_{\rm comp}$. Here $J_{\rm comp}$ is bounded independently
of the field and bit lengths. The walk preserves the desired product
measure.

Writing its gap as $\gamma>0$, we have
$\|K^d-\Pi\|_{2\to2}\le(1-\gamma)^d$ and
$\|K^d-\Pi\|_{\infty\to\infty}\le2$.
The diagonal case of the Riesz--Thorin interpolation theorem, already
covered by Riesz~\cite[Theorems~II and~II$'$, p.~472]{Riesz1927}, gives
\[
\|K^d-\Pi\|_{s\to s}
\le 2^{1-2/s}(1-\gamma)^{2d/s},
\]
which has the required form. For Thorin's complex extension, see
\cite[Introduction, \S0.1; Russian translation, pp.~43--45]{Thorin1948}.
\end{proof}

\subsection{Describe a walk before the queried vertex is available}

The previous proof describes a walk in coordinates that depend on the conditioned identifier. Saving that identifier across recursive calls would be too expensive. Instead, enumerate a bounded number of possible wrap corrections for each affine step. Each guessed correction defines an invertible transformation of the full environment without using the identifier. When the algorithm reaches the identifier, it checks which guess is valid. Exactly one survives.

This construction is useful twice: it provides short computational descriptions of averaging paths, and its invertibility permits changes of variables in the error analysis.

\begin{lemma}[Validated walk variants]\label{lem:env-variants}
A length-$d$ word in the walk of Lemma~\ref{lem:env-gap} can be expanded
into a bounded-alphabet family of variants described by $O(d)$ bits
such that:
\begin{enumerate}
\item each fixed variant acts bijectively and measure-preservingly on
the full environment before any endpoint tests;
\item its transformation is independent of the indicated identifiers;
\item for each walk word and each environment in the required
conditional slice, exactly one variant passes the slice and wrap
tests at the indicated identifiers, and this variant gives the
conditional walk transformation;
\item the sum of valid variants, each with the probability of its
original walk word, is exactly the conditional average.
\end{enumerate}
The tests can be performed when their indicated vertices are visited.
In the admissible implementation range, field elements and identifiers
have $O(B)$ bits, the total auxiliary bit string has $O(B)$ bits, and $d=O(B)$.
In that range the transformations and tests use $O(B)$ shared scratch space, and
the retained description of a variant uses $O(d)$ bits.
\end{lemma}

\begin{proof}
The only missing information is how the affine image wraps around the retention box. Since every generator has fixed coefficients, only a bounded number of wraps can occur in one step. Consider a fixed affine generator $(M,c)$ on a channel with retention
side $L$. Enumerate the transformations
\begin{equation}\label{eq:env-wrap}
A'=MA+(c+Lz)e_1^{\mathsf T}\pmod P,
\end{equation}
where $e_1$ selects the constant column and $z$ ranges over a fixed
finite box of integer vectors. That box contains all possible wraps:
for $w\in\{0,\ldots,L-1\}^H$, every coordinate of $Mw+c$ is bounded
in absolute value by a fixed multiple of $L$ plus a fixed constant.
Since $L\ge1$ and the generators are fixed, one bound for $z$ works
for every side length.

Each transformation in Equation~\eqref{eq:env-wrap} is invertible.
Left multiplication preserves the frame condition, and the added
term changes only the constant column. It therefore preserves uniform
measure on the full channel. It does not refer to the indicated
identifier.

At that identifier, let $w=Av(y)$ in its integer representatives.
Check first that $w$ is in the retention box, and then that
$Mw+c+Lz$ lies in the same box coordinatewise, using integer arithmetic
before reduction modulo $P$. The unique valid wrap has coordinates
$z_i=-\lfloor(Mw+c)_i/L\rfloor$. For this vector the evaluation of
$A'$ is the representative of the intended affine action modulo $L$.
Applying these checks at every step gives the unique valid wrap
sequence for a whole word.

The construction applies separately to the bounded number of channels.
On a channel of side $P$ the ordinary field reduction suffices;
alternatively, the same wrap validation may be used at its public
point. The free-bit torus has ordinary invertible affine operations
and requires no unknown endpoint. Frozen bits remain fixed.
Thus a word has only constantly many generator and wrap choices
per step, giving the $O(d)$-bit description.

The weights are important. These guesses are alternative descriptions of one walk step, not extra random choices that should be averaged uniformly. For each word and initial environment in the slice, exactly one variant survives. Its weight must be the weight of that word.
In particular, the sum is not divided by the number of guessed
wraps. This proves the asserted exact identity.
All matrices have fixed dimensions and $O(B)$-bit entries;
evaluating a bounded-degree polynomial or a word's individual steps
needs $O(B)$ arithmetic scratch. A test depending on a vertex can
be delayed until that vertex is current, with the required matrix
state regenerated from the recorded transformations.
Theorem~\ref{thm:ctrl-realization} accounts for reversible restoration
and the simultaneous space used when these operations are nested.
\end{proof}

\begin{remark}[Effective environment coordinates]
\label{rem:env-effective}
Conditional averaging is expressed in the channel and bit coordinates
of the environment supplied to that invocation. A fixed variant from
Lemma~\ref{lem:env-variants} is a measure-preserving change of variables
on the full environment; its retention and wrap tests remain explicit masks.
Thus free and frozen bits always refer to the supplied coordinates,
which may differ from those before the change of variables.
\end{remark}

\begin{remark}\label{rem:env-averaging}
Two distinct uses of the variants should be kept separate.
For centered-noise estimates, one first recognizes the sum of valid
variants as the conditional walk and applies
Lemma~\ref{lem:env-gap}. For crude bounds and changes of variables,
each fixed variant is a measure-preserving bijection, and one may
then pay for their $2^{O(d)}$ number. Counting variants before a
spectral estimate would needlessly weaken the useful contraction.
\end{remark}

